E-local (local features) vs E-global (global arrangement). With per-image random tile shuffling, applied in the same way to training and test images: t=16: CNN 79.1\% (R=0.75), MLP 40.2\%, t=8: CNN 66.7\% (R=0.41), MLP 32.1\%, t=4: CNN 56.6\% (R=0.12), MLP 29.0\%; clean CNN 87.8\%, permuted CNN 52.3\%. At t=8 the CNN retains R = 0.41 of its spatial benefit, and the CNN-MLP gap is +34.6 pp (clean +32.7 pp). Confidence: the CNN's correct predictions have mean confidence 0.973 (clean), 0.938 (tile8) and 0.919 (perm), so Rc = 0.36. The AUROC of confidence for separating right from wrong is 0.894 / 0.776 / 0.722, and the ECE is 0.074 / 0.230 / 0.346. The classes that lose most for the CNN at t=8 are deer, cat, bird. The MLP changes by -23.0 pp at t=8 although it was unaffected by the fixed permutation: its features depend on content appearing at consistent absolute positions. Decision (pre-registered rule R(8) $\geq$ 0.5 not met): E-GLOBAL is supported. Randomising the arrangement of 8-px tiles removes most of the CNN's advantage, so it depends on how local parts are arranged across the image. The confidence check agrees.
